\documentclass{ecscw2007}
\usepackage[T1]{fontenc}
\usepackage{ucs}
\usepackage{graphicx}
\usepackage{url}
%\usepackage{harvard}
\urlstyle{rm}

% Basically to cite any reference,
%    use \citet{Bowers91} to get 'Bowers (1991)'
%    or \cite{Bowers} to get '(Bowers, 1991)'

\clubpenalty=10000
\widowpenalty=10000
\tolerance=1
\emergencystretch=\maxdimen
\hyphenpenalty=10000
\hbadness=10000

\pagestyle{empty}

\begin{document}

%Papers Title goes here
\title{ECSCW 2024 Papers Instructions}

%Author Information goes here (Leave Empty for first Submission)
\author{Author(s)}
\affiliation{Affiliation(s)}
\email{Contact Author: Email address(es)}

%Conference Information goes here
\conferenceYear{2024}
\conferenceName{22nd European Conference on Computer-Supported Cooperative Work: The International Venue on Practice-centered Computing on the Design of Cooperation Technologies}
\paperType{Exploratory Papers} %<Exploratory Papers/Notes/Posters>
\journalName{Reports of the European Society for Socially Embedded Technologies}
\issn{XXX-XXXX}
\doi{10.18420/ecscw2024-to-be-added}

\maketitle

\begin{abstract}
This document is formatted in accordance with the ECSCW 2024 formatting guidelines for exploratory papers by means of
\LaTeX{} style sheet. You may want to use the file as a template.
\end{abstract}

\section*{Document layout}

The text bloc (the printed area) should measure and must not exceed 140 x 227 mm; it must have a one-column format and be positioned at the center of the page. In A4 format, this corresponds to 35 mm margins on all sides. If you use other formats, e.g., US Letter, please make sure the text bloc measures 140 x 227 mm (i.e., make sure line lengths do not exceed 140 mm).

No `live links' are allowed in the file. That is, all automatic links, codes etc. in the document must be switched off (e.g., copy editing tracks, hyperlinks, cross reference links, citation codes).

\section*{Title, author, abstract}

The title of the article must be preceded by 50 mm blank space (the blank space on top of the first page will be used for editorial purposes). The title should be set in 24 point Times, flushed left.

The title should be followed by the name of the author(s) in 14 point Times, flushed left. Immediately after the name of the author, the name and country of the affiliation of the author should be given in a separate line (12 point Times, flushed left).

The title, author's name, and author's affiliation should be followed by a short \textbf{abstract} of not more than 200 words which clearly summarizes the paper. The abstract should be set in 10 point Helvetica (or an equivalent sans serif font) with 13 point line leading (line space). The abstract should have 24 point space before and after.

\section*{Body text}

Body text should be set in 12 point Times (the font used here) or an equivalent serif font. Line spacing (leading) should be 15 point. The first paragraph after a headline should not be indented; otherwise the first line of a body text paragraph should be indented by 5 mm. Body text should be justified (i.e., do not use a ragged-right style).

Headings should be set in Times and flush left:
\begin{itemize}
\item Level 1: 18 point, 24 point space before, 12 point after.
\item Level 2: 14 point, 18 point space before, 9 point after.
\item Level 3: 12 point, 12 point space before, 6 point after.
\end{itemize}

\begin{quotation}
Quoted passages of more than 40 words should be set off from the text by indenting the left-hand margin 5 mm. Such quotations should be set in 10 point Times with 13 point line spacing, with 2 mm additional spacing before and after the quotation.
\end{quotation}

Trademarks and system names should be set in small caps, for example, \textsc{The Coordinator}.


\section*{References}

The reference format is based on the Springer version of the APA format, but with additional requirements:
\begin{itemize}
\item To avoid mistakes, volume numbers, issue numbers, and page numbers must be indicated by standard labels (`vol.', `no.', and `pp.', respectively).
\item Full article or book titles are required, also for proceedings.
\item To avoid misidentification, authors' names should include the first (or given) name.
\end{itemize}

Citations should be incorporated into the text, either directly in the sentence (`As claimed by \citet{Bowers91}...') or at the end, with author's name and date of publication in parenthesis: \cite{Bowers91}.

References should be listed at the end of the article, after the Notes, in alphabetical and chronological order. References should be set in 10 point Times, with 5 mm hanging indents.

(a) References to books should include the author's name; year of publication; title in full, in italics; publisher; place of publication; volume number (if appropriate).

%Johansen, R. (1988): Groupware. Computer Support for Business Teams, The Free Press, New York and London.

\begin{list}{}{\settowidth\labelwidth{}%
  \leftmargin\labelwidth
  \advance\leftmargin\labelsep}
\small
\bibitem[\protect\citeauthoryear{Johansen}{1988}]{Johansen88}
Johansen, R. (1988): {\em Groupware. Computer Support for Business Teams}.
\newblock New York and London: The Free Press.
\end{list}

(b) References to articles in periodicals should include the author's name; year of publication; title of article, in quotation marks; full title of periodical, in italics; volume number; issue number; issue date; page numbers.

%Gerson, E. M. and Star, S. L. (1986): ?Analyzing due process in the workplace', ACM Transactions on Office Information Systems, vol. 4, no. 3, July 1986, pp. 257-270.

\begin{list}{}{\settowidth\labelwidth{}%
  \leftmargin\labelwidth
  \advance\leftmargin\labelsep}
\small
\bibitem[\protect\citeauthoryear{Gerson and Star}{1986}]{Gerson86}
Gerson, E.~M. and S.~L. Star (1986): `Analyzing due process in the workplace'.
\newblock {\em ACM Transactions on Office Information Systems}, vol.~4, no.~3,
  pp. 257--270.
\end{list}

(c) References to essays in edited editions should include the author's name; year of publication; title of essay, in quotation marks; editor's name; title of collection, in italics; publisher; place of publication; volume number (if appropriate).

%Bowers, J. M. (1991): ?The Janus Faces of Design: Some Critical Questions for CSCW', in J. M. Bowers and S. D. Benford (eds.): Studies in Computer Supported Cooperative Work. Theory, Practice and Design, North-Holland, Amsterdam etc., 1991, pp. 333-350.

\begin{list}{}{\settowidth\labelwidth{}%
  \leftmargin\labelwidth
  \advance\leftmargin\labelsep}
\small
\bibitem[\protect\citeauthoryear{Bowers}{1991}]{Bowers91}
Bowers, J.~M. (1991): `{The Janus Faces of Design: Some Critical Questions for
  CSCW}'.
\newblock In: J.~M. Bowers and S.~D. Benford (eds.): {\em {Studies in Computer
  Supported Cooperative Work: Theory, Practice and Design}}. Amsterdam, etc.,
  pp. 333--350, North-Holland.
\end{list}

\section*{Notes}

Please keep notes to a minimum, and avoid long or discursive notes.  Short notes should preferably be placed as footnotes at the bottom of the page.\footnote{Footnotes should be set in 9 point Times (hanging indent 8 mm, 10 point leading). Please remember to insert a tab after the footnote mark and the footnote text.} Notes should be numbered consecutively with Arabic numerals, and indicated in the text by superscript numerals.

\section*{Tables}

Tables should be numbered consecutively with roman numbers, followed by the title (Table I. xxxxx). Horizontal rules should be indicated; vertical rules should not be used. Each table should be mentioned in the text. Each table should be positioned in place in the manuscript.

%\begin{table}[h]
%\begin{center}
%\begin{tabular}{c|c}
%x  & y \\
%w & z \\
%\end{tabular}
%\end{center}
%\label{tab:table1}
%\caption{Exemple of table.}
%\end{table}


\section*{Figures}

Figures should be numbered consecutively with arabic numbers, followed by the title (Figure 1. xxxxx). Each figure should be mentioned in the text. Each figure should be positioned in place in the manuscript:
\begin{figure}[htb]
\centerline{\includegraphics{fig1}}
\label{fig:figure1}
\caption{Figure caption in 10 point Times, 13 point leading, 12 point space before and after.}
\end{figure}

Do not use absolute picture anchors or similar gimmicks, as the result may change with different printers (\textsc{Word} is not a DTP application).

Photographs should be black-and-white and should be prepared for offset printing without further processing.

For illustrations in general, do not use proprietary graphics formats.

\section*{Long quotes}

\begin{quotation}
Long quotes should be set in 10 point Times with a 5 mm paragraph indentation and 2 mm space above and after  (as here).
\end{quotation}

\section*{Lists}

The following list formats are recommended:

\begin{itemize}
\setlength{\itemsep}{0pt}
\item list 1
\item list 1
\item list 1
\end{itemize}

\begin{enumerate}
\item list numbered 1
\item list numbered 1
\item list numbered 1
\end{enumerate}

\section*{Acknowledgments}

\small
Acknowledgment.
\normalsize

%\section*{References}
\small
\nocite{*}
\bibliographystyle{ecscw2007}
\bibliography{ecscw07_template}

\normalsize

\section*{Permissions}

If the paper is accepted for publication, a separate transfer-of-copyright form must be signed and submitted as a PDF or TIFF file, by ordinary mail, or by fax no later than the deadline for submission of camera-ready manuscripts. It is the obligation of the author(s) to obtain written permission for quotations from unpublished material, for all quotations in excess of 250 words in one extract or 500 words in total from any work still in copyright, and for the reprinting of illustrations or tables from unpublished or copyrighted material.



\end{document}
